\paragraph{C.3n (Proved): ordinary degree-one injectivity with a coefficient field, in every characteristic.}
Let $k$ be an algebraically closed field, and let
$(R,\mathfrak m,k)$ be an excellent normal strictly
henselian local $k$-algebra whose residue map restricts
to the identity on the specified coefficient field $k$.
Assume $N=\dim R\ge3$. Put
\[
 X=\operatorname{Spec}R\setminus\{\mathfrak m\},
 \qquad Y=X\cap V(f),\qquad 0\ne f\in\mathfrak m,
\]
and assume $X$ and $Y$ regular. Ordinary algebraic
equivalence is defined by actual line bundles on
the fixed $k$-scheme times a smooth connected
$k$-curve. Then
\begin{equation}\label{eq:C3-general-ordinary-injectivity}
 \ker\bigl(\operatorname{Pic}(X)_{\mathrm{alg}}
             \longrightarrow
            \operatorname{Pic}(Y)_{\mathrm{alg}}\bigr)
 \begin{cases}
   =0,&\operatorname{char}k=0,\\
   \text{is $p$-primary torsion},&\operatorname{char}k=p>0.
 \end{cases}
\end{equation}
In particular $J^1(X)\to J^1(Y)$ is injective with
rational coefficients in every characteristic;
in characteristic $p$ the integral map also becomes
injective after tensoring with $\mathbb Z[1/p]$.
Neither algebraizability nor a resolution of
singularities is assumed. The coefficient-field
hypothesis is retained explicitly; this statement
does not reinterpret ordinary algebraic equivalence
over an unspecified mixed-characteristic base.

\begin{proof}
\emph{1. An actual abelian-parameter family without
compactifying the puncture.}
We first prove the required family lemma for any
connected regular noetherian $k$-scheme $U$.
Let $T$ be a smooth connected $k$-curve with two
marked points $t_0,t_1$, and let $\mathcal M$ be
a line bundle on $U\times T$. Its smooth projective
completion is a curve $C$, with the same marked
points, and $c_0=t_0$. The product $U\times C$ is
regular and noetherian: its projection to $U$ is
smooth and of finite type. It is integral, since
$U$ is connected and regular and $C$ is
geometrically integral. Closing the divisor of a
rational section of $\mathcal M$ therefore extends
it to an invertible sheaf $\overline{\mathcal M}$
on $U\times C$. Local factoriality makes this
Weil-divisor closure Cartier.

The degree of its restriction to each curve fibre
over $U$ is locally constant. For example, this
follows from constancy of the Hilbert polynomial
for the invertible sheaf in the proper flat family
$C\times U\to U$. Since $U$ is connected, it is
one integer $d$. After the two twists
\[
 \mathcal M'=
 \overline{\mathcal M}\otimes
 \operatorname{pr}_U^*
       (\overline{\mathcal M}_{c_0}^{-1})\otimes
 \operatorname{pr}_C^*\mathcal O_C(-d c_0),
\]
the degree is zero and the restriction to
$U\times\{c_0\}$ is trivial. The last assertion
includes a choice of basis of the fixed
one-dimensional $k$-vector space
$\mathcal O_C(-dc_0)|_{c_0}$; it imposes no
condition on $U$ and requires no $k$-point of $U$.

Write $J=\operatorname{Pic}^0(C)$, and let
$\mathcal P_C$ be the Poincare bundle on $C\times J$
normalized along $c_0\times J$ and $C\times0$.
The Picard universal property gives a morphism
\[
                         g:U\longrightarrow J
\]
whose pullback of $\mathcal P_C$, with the factors
exchanged, is $\mathcal M'$.
The relevant form is
\citet[Theorem 1.1 and Lemma 1.3, pp.~2--3]{D_MilneJacobians2021}:
the point $c_0$ makes the Picard functor represented,
and normalization at $c_0$ removes the possible
line-bundle factor pulled back from the base.
Although those notes conventionally discuss
finite-type schemes, the same universal property
applies here. On each affine open of $U$, the
finite presentation of a line bundle and its
trivialization on the projective curve descend
to a finitely generated $k$-subalgebra. The
finite-type universal property gives the map
there; its uniqueness makes these maps glue.
Equivalently, this is the usual extension of
the represented Picard functor to arbitrary
$k$-schemes by affine limits.

Let $P=J^\vee$ and let $\mathcal P_J$ be the
normalized Poincare bundle on $J\times P$.
These actual universal line bundles and their
tensor-product group law are described in
\citet[Sections 9--10, especially the universal
property on p.~16]{D_MilneAbelian2022}.
Transposing $\mathcal P_C$ gives a family of
algebraically trivial line bundles on $J$
parametrized by $C$, trivial at $c_0$. It therefore
defines a morphism $h:C\to P$ with $h(c_0)=0$.
Now set
\[
             \mathcal N=(g\times1_P)^*\mathcal P_J
                    \quad\text{on }U\times P.
\]
This is an actual line bundle, trivial at $0$,
and its fibres satisfy
\begin{equation}\label{eq:C3-dual-Jacobian-family}
       [\mathcal N_{h(c)}]
        =[\overline{\mathcal M}_{c}\otimes
                     \overline{\mathcal M}_{c_0}^{-1}]
                    \quad(c\in C(k)).
\end{equation}
The degree-normalizing factor contributes only
a one-dimensional constant vector space to each
such fibre and hence does not change its line-bundle
class on $U$. This argument uses no autoduality
identification or choice of a principal polarization.
The Poincare tensor law shows that the family
$\mathcal N$ induces a homomorphism to any Picard
functor to which its line bundles map. Products
of the abelian varieties and tensor products
of these families handle finite sums of relations.

\emph{2. The exact general local Picard input.}
Apply this lemma to $U=X$. We use Boutot's local
Picard functor $G=\mathbf{Pic}^{\mathrm{loc}}_{R/k}$,
defined by sheafifying, on $k$-algebras $A$, the
Picard groups of the puncture of the henselization
of the pair
\[
                  (R\otimes_k A,\mathfrak m(R\otimes_k A)).
\]
The hypotheses and definition in
\citet[Introduction, Theorem 1(i)--(ii) and the
following paragraph, p.~VII]{D_Boutot1978}
start with an arbitrary noetherian local
$k$-algebra with residue field $k$.
They have no finite-type-over-$k$ requirement.
Depth at least two gives a finitely presented
closed unit section. Normal completion of
dimension at least three gives the additional
finite-dimensional first punctured coherent
cohomology and hence representability by a
group scheme locally of finite type over $k$.
Our $\widehat R$ is normal by excellence, so
all these hypotheses hold. The closed unit
section makes this group scheme separated.
Since $R$ is already henselian and $k$ is
algebraically closed, its $k$-points are exactly
\[
                         G(k)=\operatorname{Pic}(X).
\]
Indeed the pair henselization at $A=k$ is $R$,
and etale sheafification over the algebraically
closed field adds no new point classes.

An actual family on $X\times P$ restricts
functorially to these henselized products.
The family from step 1 therefore gives a
homomorphism $P\to G$. Its scheme-theoretic
image $B$ is a proper connected reduced
algebraic subgroup of $G$. To check the
finiteness and properness here, the morphism
from the proper finite-type scheme $P$ to
the separated scheme $G$ is proper, its
quasicompact image lies in a finite-type
open of the locally finite-type scheme $G$,
and its closed image is proper. The image
is reduced because $P$ is reduced. Over
the perfect field $k$ a reduced algebraic
group is smooth, so this proper connected
group is an abelian variety. Its inclusion
in $G$ is a closed immersion.
Thus every ordinary algebraically trivial
class on $X$ belongs to $B(k)$ for some
abelian subvariety $B\subset G$.

\emph{3. Closed restriction kernel and the
prime-to-$p$ obstruction.}
As proved in C.3k(i), $R/(f)$ is reduced,
its finite normalization is a product
\[
                            \prod_a C_a
\]
of normal strictly henselian local domains
of dimension $N-1\ge2$, and the disjoint
union of their punctures is $Y$.
Each residue field is $k$, since $k$ is
algebraically closed, and each $C_a$ has
the inherited coefficient field $k$.
Boutot's theorem therefore gives a finitely
presented closed unit section for each
functor $H_a=\mathbf{Pic}^{\mathrm{loc}}_{C_a/k}$,
even when that whole functor is not
represented by a finite-type scheme.

Restriction induces a morphism of group
functors
\[
                  G\longrightarrow\prod_a H_a.
\]
The maps on henselized tensor products
exist by the universal property of
henselization. In each finite local
factor $C_a$, the ideals
$\mathfrak m C_a$ and $\mathfrak m_{C_a}$
differ by a nilpotent ideal in the
quotient. This remains true after
tensoring with any $k$-algebra, so they
give the same henselization of the pair.
Pullback of the closed
unit section along
$B\to G\to\prod_a H_a$ defines a closed
subgroup scheme $K\subset B$. Its
$k$-points are precisely those classes
in $B(k)$ which restrict trivially to $Y$.
No universal line bundle on $B$ is used
in defining this subgroup.

For every prime $\ell$ different from
$\operatorname{char}k$, C.3k(i) gives
\[
                             K(k)[\ell]=0.
\]
If $K$ had positive dimension, its
reduced identity component would be
a positive-dimensional abelian variety.
Such a variety has nontrivial
$\ell$-torsion for every such $\ell$,
a contradiction. Hence $K$ is finite.
In characteristic zero its finite group
of $k$-points has no torsion of any
prime order and is therefore trivial.
In characteristic $p$ the finite
abelian group $K(k)$ is a $p$-group.
Every ordinary kernel element lies in
one of these groups $K(k)$, proving
\eqref{eq:C3-general-ordinary-injectivity}.
No uniform exponent for the entire
ordinary kernel is asserted.
\end{proof}

This strengthens the positive statement C.3k(ii)
in both its algebraizability and characteristic
restrictions. It is compatible with C.3m, which
refutes surjectivity, and with C.3i, whose full
Picard kernel lies outside ordinary algebraic
equivalence.
